% Banco de exercícios — capítulo 2
% O capítulo é determinado pelo nome deste ficheiro.

\begin{exo}[Do gelo frio à água em ebulição: ordens de grandeza]{chauffage-glace-eau-ebullition}
\seoready{true}
\keywords{calorimetria, mudança de estado, fusão, vaporização, calor latente, ordem de grandeza}

\textbf{Dados numéricos (à pressão atmosférica):}
\[
\begin{aligned}
c_{\text{gelo}} &= 2{,}1\times10^3\,\text{J}\!\cdot\!\text{kg}^{-1}\!\cdot\!\text{K}^{-1},
& L_f &= 3{,}34\times10^5\,\text{J}\!\cdot\!\text{kg}^{-1},\\
c_{\text{water}} &= 4{,}18\times10^3\,\text{J}\!\cdot\!\text{kg}^{-1}\!\cdot\!\text{K}^{-1},
& L_v &= 2{,}26\times10^6\,\text{J}\!\cdot\!\text{kg}^{-1}.
\end{aligned}
\]
À pressão atmosférica, aquece-se progressivamente um quilograma de gelo, inicialmente a $-100\,{}^\circ\text{C}$.

\begin{enumerate}
  \item Calcular a energia necessária para levar o gelo de $-100\,{}^\circ\text{C}$ a $0\,{}^\circ\text{C}$.
  \item Calcular a energia necessária para fundir completamente o gelo a $0\,{}^\circ\text{C}$.
  \item Calcular a energia necessária para aquecer a água líquida de $0\,{}^\circ\text{C}$ a $100\,{}^\circ\text{C}$.
  \item Calcular a energia adicional necessária para vaporizar completamente a água a $100\,{}^\circ\text{C}$.
  \item Qual das etapas requer mais energia? Comparar a energia necessária para aquecer um quilograma de água de $0\,{}^\circ\text{C}$ a $100\,{}^\circ\text{C}$ com a energia potencial de uma massa $m$ que cai de uma altura de dez metros. Que massa $m$ deve cair para igualar essa energia?
  \item Considere agora o processo inverso. Uma massa $m_v$ de vapor de água condensa-se num para-brisas, sem que a água líquida obtida arrefeça em seguida. Exprimir o calor $Q_{\text{lib}}$ libertado durante a condensação e calculá-lo para $m_v=1{,}0\times10^{-2}\,\text{kg}$. À temperatura do para-brisas, tome $L_v=2{,}45\times10^6\,\text{J}\!\cdot\!\text{kg}^{-1}$.
\end{enumerate}

\begin{indication}
Para aquecer uma fase sem mudança de estado, utilizar $Q=mc\Delta T$. Para uma mudança de estado a temperatura constante, utilizar $Q=mL$. A energia potencial gravítica de uma massa $m$ situada a uma altura $h$ é $E_p=mgh$; tome $g=9{,}81\,\text{m}\!\cdot\!\text{s}^{-2}$.
\end{indication}

\begin{solution}
\textbf{1.} O aquecimento do gelo até à sua temperatura de fusão requer
\[
Q_1=mc_{\text{gelo}}\Delta T
=1{,}0\times2{,}1\times10^3\times100
=2{,}1\times10^5\,\text{J}.
\]

\textbf{2.} Durante a fusão, a temperatura permanece igual a $0\,{}^\circ\text{C}$:
\[
Q_2=mL_f=1{,}0\times3{,}34\times10^5
=3{,}34\times10^5\,\text{J}.
\]

\textbf{3.} Para aquecer a água líquida de $0$ a $100\,{}^\circ\text{C}$,
\[
Q_3=mc_{\text{water}}\Delta T
=1{,}0\times4{,}18\times10^3\times100
=4{,}18\times10^5\,\text{J}.
\]

\textbf{4.} A vaporização completa requer ainda
\[
Q_4=mL_v=1{,}0\times2{,}26\times10^6
=2{,}26\times10^6\,\text{J}.
\]
Só esta etapa requer, portanto, uma energia da ordem de alguns megajoules.

\textbf{5.} A vaporização é, de longe, a etapa mais dispendiosa em energia:
\[
Q_4=2{,}26\times10^6\,\text{J}
>Q_3=4{,}18\times10^5\,\text{J}
>Q_2=3{,}34\times10^5\,\text{J}
>Q_1=2{,}1\times10^5\,\text{J}.
\]
Em particular, a vaporização exige cerca de $(2{,}26\times10^6)/(4{,}18\times10^5)\approx5{,}4$ vezes mais energia do que o aquecimento da água líquida de $0$ a $100\,{}^\circ\text{C}$.

Para uma massa $m$ que cai de uma altura $h=10\,\text{m}$, $E_p=mgh$. Impondo $mgh=Q_3$, obtém-se
\[
m=\frac{Q_3}{gh}
=\frac{4{,}18\times10^5}{9{,}81\times10}
\approx4{,}26\times10^3\,\text{kg}.
\]
Seria necessário fazer cair de dez metros uma massa de cerca de $4{,}3$ toneladas para libertar tanta energia como a necessária para aquecer um quilograma de água de $0$ a $100\,{}^\circ\text{C}$! É enorme e explica por que razão o aquecimento da água representa uma parte importante do consumo energético doméstico. A água possui também uma capacidade térmica mássica muito elevada, sobretudo em comparação com os metais correntes: é, por exemplo, cerca de dez vezes superior à do cobre (ver os exercícios seguintes).

\textbf{6.} A condensação é o processo inverso da vaporização: o vapor cede ao para-brisas o calor latente
\[
Q_{\text{lib}}=m_vL_v.
\]
Para $m_v=1{,}0\times10^{-2}\,\text{kg}$,
\[
Q_{\text{lib}}=1{,}0\times10^{-2}\times2{,}45\times10^6
=2{,}45\times10^4\,\text{J}.
\]
Assim, a condensação de uma pequena massa de vapor pode libertar uma quantidade de calor não desprezável. Esse calor é recebido pelo para-brisas e pelo ambiente que o rodeia.
\end{solution}
\end{exo}

\begin{exo}[A evapotranspiração de uma árvore de grande porte: um climatizador natural?]{odg-evapotranspiration-arbre}
\seoready{true}
\keywords{ordem de grandeza, evapotranspiração, árvore, calor latente, potência frigorífica, climatizador, COP}

Num dia quente e seco, uma árvore de grande porte com água suficiente pode evapotranspirar cerca de $500\,\text{L}=5{,}0\times10^{-1}\,\text{m}^3$ de água. Como a transpiração ocorre principalmente na presença de luz, supõe-se que se distribui ao longo de doze horas. A copa da árvore cobre no solo uma área $A_{\text{tree}}=50\,\text{m}^2$. São dadas a massa volúmica da água e o seu calor latente de vaporização à pressão atmosférica:
\[
\rho_{\text{water}}=1{,}0\times10^3\,\text{kg}\!\cdot\!\text{m}^{-3},
\qquad L_v=2{,}45\times10^6\,\text{J}\!\cdot\!\text{kg}^{-1}.
\]

\begin{enumerate}
  \item Calcular a massa de água evapotranspirada pela árvore durante o dia e a energia necessária para a vaporizar.
  \item Explicar por que razão este processo produz um efeito de arrefecimento.
  \item Calcular a potência média de arrefecimento da árvore por metro quadrado durante as doze horas de transpiração ativa.
  \item Para comparação, considere um climatizador alimentado por uma potência elétrica $P_{\mathrm{el}}=2{,}0\times10^3\,\text{W}$, com coeficiente de desempenho frigorífico $\mathrm{COP}=3{,}0$, que arrefece uma divisão de área $A_{\text{div}}=20\,\text{m}^2$. Por definição, $\mathrm{COP}=P_{\text{frio}}/P_{\mathrm{el}}$. Calcular a potência frigorífica e a potência frigorífica por metro quadrado, comparando-as com os resultados da árvore.
  \item Por que não se pode concluir que a árvore e o climatizador produzem exatamente a mesma sensação de arrefecimento, mesmo que as potências por unidade de área tenham a mesma ordem de grandeza?
  \item É possível arrefecer uma casa com muitas plantas de interior? (Cultura geral: a resposta depende de processos biológicos.)
\end{enumerate}

\begin{indication}
Utilizar $m=\rho V$, $Q_{\mathrm{evap}}=mL_v$ e $P=Q/\tau$. Para o climatizador, $P_{\text{frio}}=\mathrm{COP}\,P_{\mathrm{el}}$.
\end{indication}

\begin{solution}
\textbf{1.} A massa de água correspondente é
\[
m=\rho_{\text{water}}V
=1{,}0\times10^3\times5{,}0\times10^{-1}
=5{,}0\times10^2\,\text{kg}.
\]
A energia necessária para vaporizar esta água é
\[
Q_{\mathrm{evap}}=mL_v
=5{,}0\times10^2\times2{,}45\times10^6
=1{,}225\times10^9\,\text{J}.
\]
A ordem de grandeza é, portanto, um gigajoule.

\textbf{2.} A vaporização é uma mudança de estado endotérmica: requer a energia $Q_{\mathrm{evap}}$, retirada das folhas, do solo e do ar circundante. Esta extração de calor reduz a temperatura das folhas e do ambiente próximo, originando o efeito de arrefecimento da evapotranspiração.

\textbf{3.} Doze horas correspondem a
\[
\tau=12\times3600=4{,}32\times10^4\,\text{s}.
\]
A potência média de arrefecimento da árvore nesse período é
\[
P_{\text{tree}}=\frac{Q_{\mathrm{evap}}}{\tau}
=\frac{1{,}225\times10^9}{4{,}32\times10^4}
\approx2{,}84\times10^4\,\text{W}.
\]
Por unidade da área coberta pela copa,
\[
\frac{P_{\text{tree}}}{A_{\text{tree}}}
=\frac{2{,}84\times10^4}{50}
\approx5{,}7\times10^2\,\text{W}\!\cdot\!\text{m}^{-2}.
\]

\textbf{4.} A potência frigorífica útil do climatizador é
\[
P_{\text{frio}}=\mathrm{COP}\,P_{\mathrm{el}}
=3{,}0\times2{,}0\times10^3
=6{,}0\times10^3\,\text{W}.
\]
Por unidade de área da divisão,
\[
\frac{P_{\text{frio}}}{A_{\text{div}}}
=\frac{6{,}0\times10^3}{20}
=3{,}0\times10^2\,\text{W}\!\cdot\!\text{m}^{-2}.
\]
Neste modelo, a potência de evapotranspiração da árvore por unidade de área é quase duas vezes superior à de um climatizador comum.

\textbf{5.} A comparação diz respeito apenas a fluxos de energia. O climatizador arrefece um volume fechado e controlado, enquanto o vapor produzido pela árvore se dispersa na atmosfera e o vento traz ar quente. O efeito sentido depende da ventilação, da humidade, da temperatura ambiente e da distância à árvore. A transpiração varia também com a radiação solar, o vento, a humidade do ar, a espécie e a água disponível, podendo diminuir muito durante uma seca. Por outro lado, a árvore proporciona sombra, reduzindo diretamente a radiação solar recebida pelo solo e pelas pessoas. As potências calculadas permitem comparar ordens de grandeza, mas não estabelecer uma equivalência exata de conforto térmico.

\textbf{6.} Na prática, o efeito de arrefecimento das plantas de interior seria muito reduzido. Na maioria das plantas, a luz favorece a abertura dos estomas, pequenos poros das folhas por onde escapa o vapor de água. A iluminação interior é geralmente fraca: os estomas abrem menos e a transpiração diminui. Além disso, numa divisão pouco ventilada, o aumento da humidade do ar abranda progressivamente a evaporação.

Muitas plantas podem, portanto, produzir um ligeiro arrefecimento evaporativo local, mas não substituem um climatizador. Para retirar energia da habitação de modo duradouro, seria necessário renovar o ar húmido para o exterior. Uma iluminação artificial intensa poderia estimular a transpiração, mas a sua energia elétrica acabaria quase inteiramente sob a forma de calor na divisão. As plantas CAM adaptadas a ambientes áridos, incluindo os cactos, constituem uma exceção: os seus estomas abrem sobretudo à noite.
\end{solution}
\end{exo}

\begin{exo}[Mistura de duas massas de água]{calorimetrie-melange-eau}
\seoready{true}
\keywords{calorimetria, mistura, capacidade térmica, calor específico, Joseph Black, temperatura de equilíbrio}

Deitam-se $m_1=0{,}200\,\text{kg}$ de água a $T_1=80\,^\circ\text{C}$ num recipiente que já contém $m_2=0{,}300\,\text{kg}$ de água a $T_2=15\,^\circ\text{C}$. O recipiente é um calorímetro ideal: desprezam-se as perdas de calor para o exterior e a capacidade térmica do próprio recipiente. Recorda-se a relação
\[
Q=mc\Delta T,
\]
em que $c=4{,}18\times10^3\,\text{J}\!\cdot\!\text{kg}^{-1}\!\cdot\!\text{K}^{-1}$ é a capacidade térmica mássica da água líquida.

\begin{enumerate}
  \item Justificar que o calor cedido pela água quente é integralmente recebido pela água fria. Escrever a equação de conservação que envolve $m_1$, $m_2$, $c$, $T_1$, $T_2$ e a temperatura de equilíbrio $T_{eq}$.
  \item Deduzir a expressão de $T_{eq}$ e calcular o seu valor numérico.
  \item Calcular a quantidade de calor $Q$ cedida pela água quente durante a mistura.
  \item Permitem estas misturas de uma mesma substância determinar o valor de $c$? Justificar.
\end{enumerate}

\begin{indication}
Como o calorímetro é isolado e rígido, conserva-se a energia interna total $U_1+U_2$: o calor cedido por uma parte é recebido pela outra, com sinal oposto.
\end{indication}

\begin{solution}
\textbf{1.} A conservação da energia interna total fornece
\[
m_1c(T_{eq}-T_1)+m_2c(T_{eq}-T_2)=0.
\]

\textbf{2.} O fator comum $c$ simplifica-se, obtendo-se
\[
T_{eq}=\frac{m_1T_1+m_2T_2}{m_1+m_2}
=\frac{0{,}200\times80+0{,}300\times15}{0{,}200+0{,}300}
=41\,^\circ\text{C}.
\]

\textbf{3.} O calor cedido pela água quente é
\[
Q=m_1c(T_1-T_{eq})
=0{,}200\times4{,}18\times10^3\times(80-41)
\approx3{,}26\times10^4\,\text{J}.
\]
A água fria recebe exatamente a mesma quantidade:
\[
m_2c(T_{eq}-T_2)
=0{,}300\times4{,}18\times10^3\times26
\approx3{,}26\times10^4\,\text{J}.
\]

\textbf{4.} Não. Para duas massas da mesma substância, o mesmo coeficiente $c$ multiplica os dois termos do balanço e simplifica-se. A temperatura de equilíbrio medida não depende, portanto, de $c$: é apenas a média das temperaturas iniciais ponderada pelas massas. Para medir uma capacidade térmica mássica pelo método das misturas, é necessário introduzir uma substância de capacidade conhecida, como no exercício seguinte.
\end{solution}
\end{exo}

\begin{exo}[Determinação da capacidade térmica mássica de um metal pelo método das misturas]{calorimetrie-methode-melanges}
\seoready{true}
\keywords{calorimetria, método das misturas, capacidade térmica mássica, calorímetro, equivalente em água}

Pretende-se medir, como já fazia Joseph Black no século XVIII, a capacidade térmica mássica $c_m$ de um metal desconhecido pelo \emph{método das misturas}. Uma amostra de massa $m=0{,}150\,\text{kg}$ é aquecida até $T_m=95\,^\circ\text{C}$ e mergulhada rapidamente num calorímetro que contém $m_e=0{,}250\,\text{kg}$ de água, com $c_e=4{,}18\times10^3\,\text{J}\!\cdot\!\text{kg}^{-1}\!\cdot\!\text{K}^{-1}$, inicialmente a $T_e=18\,^\circ\text{C}$. Mede-se a temperatura de equilíbrio $T_{eq}=22\,^\circ\text{C}$. Numa primeira etapa, despreza-se a capacidade térmica do recipiente, do agitador e do termómetro.

\begin{enumerate}
  \item Escrever a conservação da energia do sistema isolado formado pelo metal e pela água, deixando $c_m$ como única incógnita.
  \item Deduzir a expressão de $c_m$ e calcular o seu valor. A que metal corrente corresponde aproximadamente? Valores de referência em $\text{J}\!\cdot\!\text{kg}^{-1}\!\cdot\!\text{K}^{-1}$: alumínio $9{,}0\times10^2$, ferro $4{,}5\times10^2$, cobre $3{,}9\times10^2$, chumbo $1{,}3\times10^2$.
  \item O recipiente absorve, na realidade, parte do calor. Modela-se este efeito por um \emph{equivalente em água} $\mu=1{,}5\times10^{-2}\,\text{kg}$: o calorímetro comporta-se como uma massa adicional $\mu$ de água. Recalcular $c_m$ e comentar o sentido da correção.
  \item Por que é essencial mergulhar rapidamente a amostra e isolar bem o calorímetro do ar ambiente durante a medição?
\end{enumerate}

\begin{indication}
O metal cede $Q_m=mc_m(T_m-T_{eq})$, integralmente recebido pela água e, na questão 3, também pelo calorímetro: $Q_m=(m_e+\mu)c_e(T_{eq}-T_e)$.
\end{indication}

\begin{solution}
\textbf{1.} Para o sistema isolado,
\[
mc_m(T_m-T_{eq})=m_ec_e(T_{eq}-T_e).
\]

\textbf{2.} Isolando $c_m$,
\[
c_m=\frac{m_ec_e(T_{eq}-T_e)}{m(T_m-T_{eq})}
=\frac{0{,}250\times4{,}18\times10^3\times(22-18)}{0{,}150\times(95-22)}
\approx3{,}82\times10^2\,\text{J}\!\cdot\!\text{kg}^{-1}\!\cdot\!\text{K}^{-1}.
\]
Este valor é muito próximo do valor do cobre.

\textbf{3.} Incluindo o equivalente em água do calorímetro,
\[
c_m=\frac{(m_e+\mu)c_e(T_{eq}-T_e)}{m(T_m-T_{eq})}
=\frac{(0{,}250+0{,}015)\times4{,}18\times10^3\times4}{0{,}150\times73}
\approx4{,}05\times10^2\,\text{J}\!\cdot\!\text{kg}^{-1}\!\cdot\!\text{K}^{-1}.
\]
A correção aumenta $c_m$: ao desprezar o recipiente, subestimava-se o calor realmente cedido pelo metal e, portanto, a sua capacidade térmica mássica.

\textbf{4.} O método pressupõe um sistema isolado durante toda a medição. Mergulhar rapidamente a amostra limita as trocas de calor com o ar antes de ela chegar à água; um bom isolamento limita as perdas enquanto se estabelece o equilíbrio. Qualquer transferência não contabilizada falsearia o balanço e o valor deduzido de $c_m$. É precisamente este tipo de cuidado experimental que já era aplicado por Black e, mais tarde, por Lavoisier e Laplace com o seu calorímetro de gelo.
\end{solution}
\end{exo}

\begin{exo}[Calorias alimentares e potência metabólica]{calorimetrie-calories-alimentaires}
\seoready{true}
\keywords{caloria, quilocaloria, Caloria alimentar, equivalente mecânico, potência metabólica, unidades}

Um rótulo nutricional indica que um adulto consome em média cerca de $2000\,\text{kcal}$ por dia. A quilocaloria é uma unidade não SI ainda usada em nutrição; a sua conversão para o SI é $1\,\text{kcal}=4{,}18\times10^3\,\text{J}$.

\begin{enumerate}
  \item Converter esta ingestão diária em joules.
  \item Deduzir a potência metabólica média, em watts, supondo que a energia é utilizada uniformemente durante 24 horas.
  \item Uma barra energética indica «$210\,\text{kcal}$». Que massa de água inicialmente a $20\,^\circ\text{C}$ poderia ser levada à ebulição, a $100\,^\circ\text{C}$, com a mesma energia, se esta fosse integralmente convertida em calor? Tome $c_{\text{water}}=4{,}18\times10^3\,\text{J}\!\cdot\!\text{kg}^{-1}\!\cdot\!\text{K}^{-1}$.
  \item Sob que formas é esta energia realmente utilizada pelo organismo humano? Responder qualitativamente.
\end{enumerate}

\begin{indication}
Para a questão 2, utilizar $\mathcal P=E/\Delta t$. Para a questão 3, converter primeiro a energia em joules e depois isolar $m$ em $Q=mc\Delta T$.
\end{indication}

\begin{solution}
\textbf{1.} Em unidades SI,
\[
E=2000\times4{,}18\times10^3=8{,}36\times10^6\,\text{J}.
\]

\textbf{2.} A potência média é
\[
\mathcal P=\frac{8{,}36\times10^6}{8{,}64\times10^4}\approx97\,\text{W}.
\]
Ao longo de um dia, uma pessoa utiliza em média aproximadamente a mesma potência que uma lâmpada incandescente. Esta ordem de grandeza, muitas vezes surpreendente, ilustra a utilidade de converter as calorias alimentares em unidades físicas habituais.

\textbf{3.} A energia da barra, em joules, é
\[
Q=210\times4{,}18\times10^3\approx8{,}78\times10^5\,\text{J}.
\]
Com $\Delta T=80\,\text{K}$,
\[
m=\frac{Q}{c\Delta T}
=\frac{8{,}78\times10^5}{4{,}18\times10^3\times80}
\approx2{,}63\,\text{kg}.
\]
Uma única barra energética representa, em ordem de grandeza, energia suficiente para levar cerca de $2{,}5\,\text{kg}$ de água da torneira à ebulição.

\textbf{4.} O organismo não «consome» energia no sentido de a fazer desaparecer: transforma a energia química dos nutrientes. Uma parte serve para o trabalho muscular, o funcionamento dos órgãos, a manutenção dos gradientes elétricos das células, o transporte de moléculas e as sínteses químicas. Outra parte pode ser armazenada quimicamente sob a forma de glicogénio ou gordura. Por fim, a maior parte da energia utilizada é dissipada como calor, contribuindo para manter a temperatura corporal antes de ser transferida para o ambiente. Uma pequena fração sai também do organismo como energia química ainda contida nos resíduos.
\end{solution}
\end{exo}
